Bei einem Kreisprozess wird Luft zunächst isotherm aufs halbe Volumen komprimiert, anschliessend isobar aufs Anfangsvolumen gebracht und schliesslich wieder in den Anfangszustand zurückgeführt.
\begin{abcliste}
\abc Skizzieren Sie den Kreisprozess in einem $p(V)$-Diagramm.
\abc Leiten Sie einen algebraischen Ausdruck für die vom Gas während eines Zyklus abgegebene Arbeit her.
\abc Berechnen Sie die Arbeit für die Anfangswerte \VaO\ und \paO.
\end{abcliste}
